\documentclass[11pt,reqno]{amsart}
\usepackage[margin=1.1in]{geometry}
\usepackage{amsmath,amssymb,amsthm,mathtools,booktabs}
\usepackage[colorlinks=true,linkcolor=blue!50!black,citecolor=blue!50!black,urlcolor=blue!50!black]{hyperref}
\usepackage{xcolor}
\usepackage[expansion=false]{microtype}
\theoremstyle{plain}
\newtheorem{theorem}{Theorem}[section]
\newtheorem{proposition}[theorem]{Proposition}
\newtheorem{lemma}[theorem]{Lemma}
\newtheorem{corollary}[theorem]{Corollary}
\newtheorem{question}[theorem]{Question}
\theoremstyle{definition}
\newtheorem{definition}[theorem]{Definition}
\newtheorem{example}[theorem]{Example}
\newtheorem*{remark*}{Remark}
\newcommand{\Z}{\mathbb{Z}}
\newcommand{\Q}{\mathbb{Q}}
\newcommand{\F}{\mathbb{F}}
\newcommand{\leg}[2]{\left(\frac{#1}{#2}\right)}
\DeclareMathOperator{\Tr}{Tr}
\DeclareMathOperator{\Nm}{N}
\newcommand{\lean}[1]{\texttt{\detokenize{#1}}}
\newcommand{\leanline}[1]{\par\smallskip{\raggedright\leftskip=1.5em\noindent{\small\emph{Lean:} #1}\par}\smallskip}
\emergencystretch=3em
\newcommand{\Pstar}{\mathcal P^*}

\begin{document}

\title[The coordinates of a quartic $L$-value]{The coordinates of a quartic $L$-value:\\
a formalized survey of the classical facts}
\author{Joseph S.\ Farland}
\date{30 October 2026 (version 2)}

\begin{abstract}
For primes $p>5$ with $p\equiv5\pmod8$, the quartic generalized Bernoulli number $\beta=B_{1,\chi_4}$ is,
after a choice of quartic character, a Gaussian integer $k+it$. This survey collects in one place the
classical facts about these coordinates that are otherwise spread across papers from 1976 to 2026: why
order four is the first level at which an odd non-real character occurs at prime conductor; what the
trace and the relative class number already determine ($k$ and $|t|$); how the normalization
$\chi_4(2)=i$ fixes the sign of $t$ and agrees with the quartic Jacobi-sum normalization and with Hudson's
quarter-range coordinates; and the fixed-base digit identity together with what it does not reach. A
section on the higher $2$-power levels records, for characters of order $2^a$ modulo primes with
$v_2(p-1)=a$, that the coset differences vanish below the top level, that the coordinates of $B_{1,\chi}$
in the power basis are odd integers, and that their first coordinate is the trace of Murty and
Thangadurai. Nothing in this survey is new; every statement is attributed to its source, and the one
statement we did not find in print, the parity of the coordinates at levels $2^a\ge8$, is the argument of
the level-four case verbatim. The elementary core (Propositions 1.1, 3.1, 3.2, 3.4 except one remark, 3.5,
4.1 and 5.1 to 5.3, Definition 3.3, and the conjugation and reconstruction statements) is formally
verified in Lean~4 with Mathlib; Section~\ref{sec:lean} lists what is verified and what is quoted.
\end{abstract}
\maketitle

\section{The classical boundary}\label{sec:boundary}

In 1839 Dirichlet proved that for a prime $p\equiv3\pmod4$, $p>3$, the interval $(0,p/2)$ contains more
quadratic residues than non-residues, and that the excess is $(2-\leg2p)h(-p)$, that is, $h(-p)$ when
$p\equiv7\pmod8$ and $3h(-p)$ when $p\equiv3\pmod8$ \cite[(1.1)]{Berndt}. The programme that grew out of this
asks, for each interval and each class of primes, whether the corresponding character sum has constant
sign and what class number measures it. Berndt's 1976 survey \cite{Berndt} is the systematic treatment: it
works through intervals of length $k/3$, $k/5$, $k/6$, $k/8$, $k/10$, $k/12$, $k/16$ and $k/24$, expresses the
sums in terms of Dirichlet $L$-values at $s=1$, translates the real cases into class numbers, and deduces
congruences.

Twice in that paper Berndt identifies the obstruction. At the sixteenth range \cite[\S10]{Berndt} the
individual sums involve $L$-functions of complex characters, so his methods yield no conclusion about their
signs; the theorem that follows concerns paired sums. In the closing section \cite[\S14]{Berndt} he states
the general point: every such interval sum can be written as a linear combination of $L$-functions at
$s=1$, but when non-real characters enter the sign argument stops. At the fifth range \cite[\S5]{Berndt} he
names quartic characters as the obstruction to a positivity conclusion.

The same boundary appears in the digit literature. Girstmair \cite[Rem.~5]{GirDecOct} treats the frequency
of the digit $0$ in the decimal period of $1/p$ when $10$ has order $(p-1)/2$. For $p\equiv3\pmod4$ the answer
is a combination of $h(-p)$ and $h(-5p)$. For $p\equiv1\pmod4$ the evaluation supplied by \cite[(6)]{SUZ}
involves a first generalized Bernoulli number attached to a character of order four, and Girstmair reports
no class-number interpretation of the individual complex value.

Classical methods do not stop at this boundary. They continue to control Galois-invariant combinations,
namely traces, norms, determinants and class-number formulas. What they do not give is an ordered sign for
the individual quartic value, which does not collapse to a signed scalar class number in the quadratic
manner. What is left over is the ordered pair of coordinates after the conjugation ambiguity has been
fixed.

\begin{proposition}[the first non-real odd level]\label{prop:boundary}
Let $\chi$ be a Dirichlet character modulo a prime $p$ and suppose that $\chi$ has exact order $q>1$. For $q=2$
the character is real. For odd $q$ (in particular $q=3$) every such character is even. At $q=4$ the character
is odd and non-real precisely when $p\equiv5\pmod8$. Thus order four is the first character order at which an
odd non-real character can occur at prime conductor.
\end{proposition}

\begin{proof}
If $q$ is odd, then $\chi(-1)$ has order dividing both $q$ and $2$, hence $\chi(-1)=1$. A character of order two
takes only the values $0$ and $\pm1$. If $\chi_4$ has order four, choose a generator $g$ of $(\Z/p\Z)^\times$ with
$\chi_4(g)=i$. Since $-1=g^{(p-1)/2}$, $\chi_4(-1)=i^{(p-1)/2}=(-1)^{(p-1)/4}$. Hence $\chi_4$ is odd exactly when
$(p-1)/4$ is odd, equivalently $p\equiv5\pmod8$. In that case $\chi_4$ and $\overline{\chi_4}$ are the two odd
characters of order four and are non-real.
\end{proof}
\leanline{\lean{apply_neg_one_of_odd_orderOf}, \lean{apply_of_orderOf_two}, \lean{quartic_apply_neg_one},
\lean{quartic_apply_neg_one_of_one}.}

This elementary boundary is the one visible in Girstmair's coset-sum method. His sign-vector argument
\cite{GirAbh} requires a nontrivial quadratic character among the odd characters, a condition that holds in
the real cases but fails at index four; his 2026 survey \cite{GirSurv} accordingly treats $q=2$, $6$ and $10$
rather than $q=4$.

\section{What classical theory determines}\label{sec:classical}

Throughout, $p$ is a prime with $p>5$ and $p\equiv5\pmod8$; write $\Pstar$ for this set. Let $\chi_4$ be one of the
two primitive quartic characters modulo $p$ and set
\[
\beta=B_{1,\chi_4}=\frac1p\sum_{1\le n\le p-1}\chi_4(n)\,n.
\]
As shown in Proposition~\ref{prop:coset}, $\beta$ lies in $\Z[i]$; write $\beta=k+it$. Let $\langle x\rangle$
denote the least nonnegative residue of $x$ modulo $p$, let $H$ be the subgroup of fourth powers modulo $p$, and,
for $j=0,1,2,3$, put
\[
C_j=\{n\in(\Z/p\Z)^\times:\chi_4(n)=i^j\},\qquad S_j=\frac1p\sum_{n\in C_j}\langle n\rangle .
\]
Thus $C_0=H$. Up to the cyclic indexing determined by $\chi_4$, the $S_j$ are the classical coset sums and, in
Girstmair's terminology \cite{GirAbh}, the components of the $b$-deviation vector at $b=p+1$.

\emph{The $L$-value normalization.} For every nonprincipal character $\chi$, $L(0,\chi)=-B_{1,\chi}$. If $\chi$
is primitive and odd of conductor $p$, the functional equation gives
\[
L(1,\overline\chi)=\frac{\pi i\,\tau(\overline\chi)}p\,B_{1,\chi},
\]
where $\tau(\overline\chi)$ is the Gauss sum. Thus $B_{1,\chi}$ is the algebraic normalization of the
corresponding value at $s=1$, and the coordinates in this survey are the coordinates of that normalized quartic
$L$-value.

\subsection{Trace, norm and the relative class number}
Let $K_p$ be the imaginary cyclic quartic subfield of $\Q(\zeta_p)$, and let $h^*(p)=h(K_p)/h(\Q(\sqrt p))$ be its
relative class number.

\begin{proposition}[the classical data]\label{prop:data}
For $p\in\Pstar$,
\begin{equation}\label{eq:trnorm}
\Tr(\beta)=2k,\qquad\Nm(\beta)=k^2+t^2=2h^*(p).
\end{equation}
Consequently the trace and the relative class number determine $k$ and $|t|$, and the only ambiguity in the
individual quartic value is the conjugation $t\mapsto-t$.
\end{proposition}

\begin{proof}
The trace identity is formal from $\beta=k+it$; its occurrence as a classical digit invariant is
\cite[Thm.~1]{MT}, restated at the $2$-power levels in \cite[Thm.~1]{GirMean} (see Proposition~\ref{prop:trace}).
Also $\beta=(S_0-S_2)+i(S_1-S_3)$. Hudson and Williams \cite{HW81}, with proof and discussion in Hudson
\cite[(2.3)]{Hudson}, give $h^*(p)=\frac12\bigl((S_2-S_0)^2+(S_3-S_1)^2\bigr)$. Substitution gives
$\Nm(\beta)=2h^*(p)$. Equivalently, this is the analytic class number formula $h^*=Qw\prod(-\frac12B_{1,\chi})$
over the two odd characters, with $Q=1$ and $w=2$.
\end{proof}
\leanline{the trace, as \lean{trace_eq}. The identification of the norm with $2h^*(p)$ needs class numbers of
number fields and is quoted, not formalized.}

The trace fixes the real coordinate and the relative class number fixes the absolute value of the imaginary
coordinate, so classical theory determines the quartic value up to complex conjugation. The role of the pin
in Section~\ref{sec:pin} is to choose between the two conjugates.

The same norm structure persists in higher degree. Girstmair \cite{GirMC}, restated precisely by Ku\v{c}era
\cite{Kucera}, expresses relative class numbers of imaginary abelian fields of prime-power conductor as group
determinants in the coset sums, with explicit formulas up to degree ten. On the trace side, Murty and
Thangadurai \cite[Thm.~1]{MT} show that for $r\mid p-1$, with the relevant characters odd, a digit sum of $1/p$
is an elementary term plus a trace of first generalized Bernoulli numbers. At $r=4$ that trace is $2k$; at the
higher $2$-power levels it is Proposition~\ref{prop:trace}.

\subsection{Interval sums}
The general evaluation underlying the interval formulas is \cite[(6)]{SUZ}: for a character $\chi$ modulo
$M$, any multiple $N$ of $M$, any $r$ prime to $N$ and any $m\ge1$,
\begin{equation}\label{eq:suz}
m\,r^{m-1}\sum\chi(n)\,n^{m-1}=-B_{m,\chi}\,r^{m-1}+\frac{\overline\chi(r)}{\varphi(r)}\sum_\psi
\overline\psi(-N)\,B_{m,\chi\psi}(N).
\end{equation}
Szmidt, Urbanowicz and Zagier trace this formula to Lerch \cite{Lerch}, Berndt \cite{Berndt} and Urbanowicz
\cite{Urb}. In \eqref{eq:suz} the first sum is over $0<n<N/r$ and the second is over characters $\psi$ modulo
$r$; $B_{m,\chi\psi}(X)=\sum_{0\le j\le m}\binom mjB_{j,\chi\psi}X^{m-j}$ is the generalized Bernoulli polynomial
of $\chi\psi$, regarded as a character modulo $rN$, evaluated at $X=N$. In \cite{SUZ} the formula is printed with
$\chi(r)$ and $\psi(-N)$. That form agrees with \eqref{eq:suz} whenever $\chi(r)$ and every $\psi$ modulo $r$ are
real, which covers every application in \cite{SUZ} (quadratic $\chi$ and $r\mid24$). For non-real characters the
conjugated form \eqref{eq:suz} is the correct one, and it is the form Girstmair applies to characters of order
four and eight \cite[(27)]{GirMean}. For instance, with $\chi=\chi_4$ normalized as in Definition~\ref{def:pin},
$p=13$, $N=p$, $r=2$ and $m=1$, the left side of \eqref{eq:suz} is $1+3i$.

At $m=1$, $N=p$ and $r=2$, formula \eqref{eq:suz} for an odd character is the fixed-base identity
\eqref{eq:fixed} at $b=2$: $\sum_{0<n<p/2}\chi(n)=(\overline\chi(2)-2)B_{1,\chi}$. At $m=2$ the formula gives the
moment-weighted forms appearing in \cite[\S13]{Berndt}. The eighth- and sixteenth-range formulas in Berndt
belong to the same linear theory as the trace above; what they do not supply is an ordering of the two
conjugate quartic values.
\leanline{the case $m=1$, $N=p$, $r=2$, as \lean{half_sum}. The general formula \eqref{eq:suz} is quoted.}

\section{The coordinates and the pin}\label{sec:pin}

\begin{proposition}[pairing and integrality]\label{prop:coset}
For $p\in\Pstar$, the four $S_j$ are integers, $S_0+S_2=S_1+S_3=(p-1)/4$, and
\begin{equation}\label{eq:kt}
k=S_0-S_2,\qquad t=S_1-S_3.
\end{equation}
\end{proposition}

\begin{proof}
Put $m=(p-1)/4$. Since $p>5$, $m>1$, so $H$ contains an element $a\ne1$; multiplication by $a$ permutes each coset
$C_j$, so the sum of its elements is $0$ in $\F_p$, and therefore the integer sum of the least positive
representatives in $C_j$ is divisible by $p$. Hence $S_j\in\Z$. Multiplication by a generator permutes the four
cosets cyclically, so each has $m$ elements. Since $p\equiv5\pmod8$, $\chi_4(-1)=-1$, so $n\mapsto p-n$
interchanges $C_j$ and $C_{j+2}$; pairing the $m$ elements gives $S_j+S_{j+2}=m$. Finally
$B_{1,\chi_4}=S_0+iS_1-S_2-iS_3$, which gives \eqref{eq:kt}.
\end{proof}
\leanline{\lean{p_dvd_cosetSum}, \lean{card_class}, \lean{S_add_S_neg}, \lean{T_eq}.}

\begin{proposition}[oddness]\label{prop:odd}
For $p\in\Pstar$, $k$ and $t$ are odd integers.
\end{proposition}

\begin{proof}
By Proposition~\ref{prop:coset}, $k=2S_0-(p-1)/4$ and $t=2S_1-(p-1)/4$. Since $(p-1)/4$ is odd, both coordinates
are odd.
\end{proof}
\leanline{\lean{odd_k}, \lean{odd_t}.}

There are exactly two primitive quartic characters modulo $p$, $\chi_4$ and $\overline{\chi_4}$. Conjugation fixes
$k$ and sends $t$ to $-t$. By Proposition~\ref{prop:data} the classical invariants already determine $k$ and $|t|$,
so the entire remaining ambiguity is the sign of the second coordinate.
\leanline{\lean{quartic_eq_or_eq_pow_three}, \lean{k_pow_three}, \lean{t_pow_three}.}

\begin{definition}[the pin]\label{def:pin}
For $p\in\Pstar$, let $\chi_4$ be the unique quartic character satisfying $\chi_4(2)=i$. This is possible and
unique because $p\equiv5\pmod8$ makes $2$ a quadratic non-residue, so $\chi_4(2)\in\{i,-i\}$, and the two quartic
characters take the two values.
\end{definition}
\leanline{\lean{quartic_apply_two}, \lean{existsUnique_pinned}.}

\begin{proposition}[the Jacobi-sum dictionary]\label{prop:jacobi}
Let $\chi_4$ be pinned by $\chi_4(2)=i$ and write $J(\chi_4,\chi_4)=a+bi$. Then $a^2+b^2=p$, $a\equiv-1\pmod4$ and
$b\equiv2\pmod4$. Replacing $\chi_4$ by $\overline{\chi_4}$ leaves $a$ fixed and sends $b$ to $-b$. Thus the pin on
$\chi_4$ and the sign choice for the second coordinate of the quartic Jacobi sum are the same conjugation choice.
\end{proposition}

\begin{proof}
This is the standard quartic Jacobi-sum evaluation \cite[Cor.~3.10]{BEW}. The norm is $|J(\chi,\chi)|^2=p$ for $\chi$
and $\chi^2$ nontrivial. For the congruences, pair $x$ with $1-x$ in $J=\sum\chi(x)\chi(1-x)$: the fixed point
$x=1/2$ contributes $\chi(1/4)=\chi(2)^{-2}=-1$, and the remaining terms pair off, so $J=2W-1$ where $W$ is a sum of
$(p-3)/2$ values in $\{\pm1,\pm i\}$. Then $a=2\operatorname{Re}W-1$ and $b=2\operatorname{Im}W$ with
$\operatorname{Re}W+\operatorname{Im}W$ odd, and $a^2+b^2=p\equiv5\pmod8$ forces $\operatorname{Im}W$ odd and
$\operatorname{Re}W$ even. Complex conjugation commutes with the finite sum defining $J$, so
$J(\overline{\chi_4},\overline{\chi_4})=a-bi$, and the pin chooses one of the two signs of $b$.
\end{proof}
\leanline{\lean{jacobiSum_norm}, \lean{jacobiSum_mod_four}, \lean{jacobiSum_pow_three}.}

\begin{remark*}[the sign selected by the pin]
If $\chi_4(2)=i$, then $b\equiv-2\pmod8$. For either quartic character $\chi$, writing $J(\chi,\chi)=a+bi$, one has
$\chi(2)=i^{ab/2}$: the Jacobi sum of the quartic residue character $\chi_\pi$ modulo a primary prime $\pi$ of norm
$p$ satisfies $-\chi_\pi(-1)J(\chi_\pi,\chi_\pi)=\pi$, and the supplementary law gives $(2/\pi)_4=i^{ab/2}$ for
$\pi=a+bi$ \cite[Ch.~9]{IR}. Here $\chi(-1)=-1$, and $a+bi$ is primary because $a\equiv3\pmod4$ and
$b\equiv2\pmod4$. Hence $\chi_4(2)=i$ exactly when $ab/2\equiv1\pmod4$; since $a\equiv3\pmod4$, this means
$b\equiv-2\pmod8$. For $p=13$, $J(\chi_4,\chi_4)=3-2i$.
\end{remark*}
\leanline{not formalized; the supplementary law for the quartic character of $2$ is quoted.}

\begin{proposition}[quarter-count dictionary]\label{prop:quarters}
Let $N_r=\#\{n\in H:rp/4<n<(r+1)p/4\}$ for $0\le r\le3$, and put $A=N_0-N_3$ and $B=N_1-N_2$. Under the
normalization $\chi_4(2)=i$, Hudson's formulas \cite[Thm.~3.1]{Hudson} give
\begin{equation}\label{eq:hudson}
\begin{aligned}
S_0&=\tfrac12\bigl(\tfrac{p-1}4-\tfrac35A-\tfrac15B\bigr),&\qquad S_1&=\tfrac12\bigl(\tfrac{p-1}4-\tfrac15A+\tfrac35B\bigr),\\
S_2&=\tfrac12\bigl(\tfrac{p-1}4+\tfrac35A+\tfrac15B\bigr),&\qquad S_3&=\tfrac12\bigl(\tfrac{p-1}4+\tfrac15A-\tfrac35B\bigr).
\end{aligned}
\end{equation}
Consequently
\begin{equation}\label{eq:ktAB}
k=-\frac{3A+B}5,\qquad t=\frac{-A+3B}5,
\end{equation}
and conversely
\begin{equation}\label{eq:ABkt}
A=-\frac{3k+t}2,\qquad B=\frac{3t-k}2.
\end{equation}
Thus $(k,t)$ and the two independent quarter-range differences are equivalent data. The four individual quarter
counts are not determined by $k$ and $t$ alone.
\end{proposition}

\begin{proof}
Equation \eqref{eq:hudson} is \cite[Thm.~3.1]{Hudson} in the normalization $\chi_4(2)=i$. A short route to it: let
$Q_r$ be the sum of $\chi_4$ over the $r$-th quarter of $(0,p)$. The fixed-base identity \eqref{eq:fixed} at $b=2$
and $b=4$, with $\chi_4(2)=i$ and $\chi_4(4)=-1$, together with the reflection $Q_{3-r}=-Q_r$, gives
$2Q_0=(-3+i)\beta$ and $2Q_1=(-1-3i)\beta$. Since $n\mapsto p-n$ carries $C_2\cap(\text{quarter }r)$ onto
$H\cap(\text{quarter }3-r)$, $\operatorname{Re}Q_0=A$ and $\operatorname{Re}Q_1=B$. Taking real parts gives
\eqref{eq:ABkt}; solving gives \eqref{eq:ktAB}; and \eqref{eq:ktAB} with Proposition~\ref{prop:coset} gives
\eqref{eq:hudson}.
\end{proof}
\leanline{\lean{two_A}, \lean{two_B} for \eqref{eq:ABkt}, \lean{five_k}, \lean{five_t} for \eqref{eq:ktAB},
\lean{hudson_S} for \eqref{eq:hudson}, cleared of denominators.}

\begin{theorem}[reconstruction]\label{thm:recon}
For $p\in\Pstar$, the classical invariants determine the quartic value up to conjugation, and the pin determines
it completely. More precisely, $k=\frac12\Tr(\beta)$ and $|t|=\sqrt{2h^*(p)-k^2}$, and $\chi_4(2)=i$ selects the
sign of $t$. With that choice, Proposition~\ref{prop:jacobi} fixes the corresponding sign of $b$ in the Jacobi sum,
and Proposition~\ref{prop:quarters} reconstructs the quarter-range differences $A$ and $B$:
\[
(\Tr\beta,\ h^*(p))\ \longrightarrow\ (k,|t|)\ \xrightarrow{\ \chi_4(2)=i\ }\ (k,t)\ \longleftrightarrow\ (A,B).
\]
\end{theorem}

\begin{proof}
The first two assertions are Proposition~\ref{prop:data} and Definition~\ref{def:pin}. The two dictionaries are
Propositions~\ref{prop:jacobi} and~\ref{prop:quarters}.
\end{proof}
\leanline{\lean{reconstruction} (every quartic $\psi$ has $k(\psi)=k(\chi_4)$ and $t(\psi)=\pm t(\chi_4)$, $\psi(2)=i$
exactly when $\psi=\chi_4$, and $(k,t)\leftrightarrow(A,B)$); the step through $h^*(p)$ is quoted.}

\begin{example}[$p=13$]\label{ex:13}
Here $H=\{1,3,9\}$, so $(N_0,N_1,N_2,N_3)=(2,0,1,0)$, hence $A=2$ and $B=-1$. Formula \eqref{eq:ktAB} gives
$k=t=-1$, so the pinned value is $\beta=-1-i$. Its norm is $2$, and Proposition~\ref{prop:data} gives
$h^*(13)=1$. The example checks the integrality, parity, norm and quarter-count conventions together.
\end{example}

\section{Fixed-base and diagonal families}\label{sec:fixed}

The classical digit and coset-sum identities use a fixed base $b$. Multiplication by $b$ is a bijection of
$(\Z/p\Z)^\times$, and composing a character with that map contributes a constant factor. For $b\ge2$ with $p\nmid b$,
put
\[
\theta(b;n)=\frac{b\langle n\rangle-\langle bn\rangle}p=\Bigl\lfloor\frac{bn}p\Bigr\rfloor,\qquad0<n<p.
\]

\begin{proposition}[fixed-base collapse]\label{prop:fixed}
For a character $\chi$ modulo $p$,
\begin{equation}\label{eq:fixed}
\sum_{1\le n\le p-1}\overline\chi(n)\,\theta(b;n)=(b-\chi(b))\,B_{1,\overline\chi}.
\end{equation}
\end{proposition}

\begin{proof}
Since $\theta(b;n)=(bn-\langle bn\rangle)/p$, the first term contributes $bB_{1,\overline\chi}$. In the second term the
substitution $m=\langle bn\rangle$ is a bijection of $(\Z/p\Z)^\times$ and $\overline\chi(b^{-1}m)=\chi(b)\overline\chi(m)$,
so it contributes $\chi(b)B_{1,\overline\chi}$. Subtraction gives \eqref{eq:fixed}.
\end{proof}
\leanline{\lean{fixed_base_collapse}, for characters with values in any commutative ring; no parity hypothesis is
needed.}

This is the prime-conductor form of \cite[Thm.~1]{GirSurv}. The mechanism is the bijective constant-factor
substitution. On the diagonal $\theta(n;n)=\lfloor n^2/p\rfloor$ the base moves with the summand: characters still
satisfy $\chi(n^\ell)=\chi(n)^\ell$, but $n\mapsto n^\ell$ is $\gcd(\ell,p-1)$-to-one onto the $\ell$-th powers, not a
bijection carrying a constant character factor, so the proof of Proposition~\ref{prop:fixed} does not apply to
diagonal sums.

\section{The higher \texorpdfstring{$2$}{2}-power levels}\label{sec:levels}

The coset description of Section~\ref{sec:pin} is the level-four case of a description at every $2$-power level.
In this section $p$ is a prime with $v_2(p-1)=a\ge2$, $g$ is a primitive root modulo $p$, and $\chi$ is a character of
exact order $2^a$, with $\zeta=\chi(g)$ a primitive $2^a$-th root of unity. Exactly as in
Proposition~\ref{prop:boundary}, $\chi(-1)=\zeta^{(p-1)/2}=-1$, because $(p-1)/2^a$ is odd: $\chi$ is odd. For $b\le a$
let $H_{2^b}$ be the subgroup of $2^b$-th powers, of order $(p-1)/2^b$; the classes of $\chi$ are the cosets
\[
C_r=\{n:\chi(n)=\zeta^r\}=g^rH_{2^a},\qquad S_r=\frac1p\sum_{n\in C_r}\langle n\rangle\qquad(0\le r<2^a),
\]
and $B_{1,\chi}=\sum_r\zeta^rS_r$. At $a=2$ these are the $C_j$ and $S_j$ of Section~\ref{sec:classical}.

\begin{proposition}[the lower levels vanish]\label{prop:even}
Let $b<a$. Every coset $C$ of $H_{2^b}$ is closed under $n\mapsto-n$, so $\frac1p\sum_{n\in C}\langle n\rangle=|C|/2$, and all
coset differences at level $2^b$ vanish. Equivalently, $B_{1,\psi}=0$ for every even character $\psi\ne1$ modulo $p$.
\end{proposition}

\begin{proof}
Since $2^b$ divides $(p-1)/2$, $-1=g^{(p-1)/2}$ lies in $H_{2^b}$. Pairing $n$ with $p-n$ in $C$ gives $2\sum_C\langle n\rangle=|C|p$.
For the character form, substitute $n\mapsto-n$ in $\sum\psi(n)\langle n\rangle$: it equals $\sum\psi(n)(p-\langle n\rangle)$, so
$2\sum\psi(n)\langle n\rangle=p\sum\psi(n)=0$. This is the classical vanishing of $B_{1,\psi}$ for even $\psi\ne1$
\cite[\S4.1]{Washington}.
\end{proof}
\leanline{\lean{levelSum_even} (coset form, for the classes of any even character), \lean{even_T_eq_zero}.}

So among the characters of order dividing $2^a$ only those of exact order $2^a$ have nonzero $B_{1,\chi}$, and only the
top level carries coordinates.

\begin{proposition}[integrality and parity at every level]\label{prop:levels}
Suppose that $p$ is not a Fermat prime, that is, $|H_{2^a}|=(p-1)/2^a>1$. Then the $S_r$ are integers,
$S_r+S_{r+2^{a-1}}=|H_{2^a}|$, and
\[
B_{1,\chi}=\sum_{0\le r<2^{a-1}}d_r\zeta^r,\qquad d_r=S_r-S_{r+2^{a-1}},
\]
with every $d_r$ an odd integer.
\end{proposition}

\begin{proof}
This is the proof of Propositions~\ref{prop:coset} and~\ref{prop:odd}. An element $c\ne1$ of $H_{2^a}$ permutes each $C_r$,
so the sum of $C_r$ vanishes in $\F_p$ and $S_r\in\Z$. Since $\chi$ is odd, $n\mapsto p-n$ carries $C_r$ onto
$C_{r+2^{a-1}}$, which gives the pairing; and $\zeta^{r+2^{a-1}}=-\zeta^r$ gives the expansion. Finally
$d_r=2S_r-|H_{2^a}|$ with $|H_{2^a}|$ odd: a subgroup of the cyclic group $(\Z/p\Z)^\times$ of even order would contain its
element of order two, which is $-1$, and $\chi(-1)=-1\ne1$.
\end{proof}
\leanline{\lean{p_dvd_levelSum}, \lean{levelSum_add_neg}, \lean{level_parity}; the oddness of $|H_{2^a}|$ is
\lean{card_kernel_odd}.}

At a Fermat prime the subgroup is trivial and integrality fails: at $p=17$, $a=4$, one has
$d_0=2\cdot\frac1{17}-1=-\frac{15}{17}$. The integrality of coset sums of a nontrivial subgroup is classical
\cite{GirAbh,GirSurv}, and so is the parity at level four (Proposition~\ref{prop:odd}); we have not found the parity at
the levels $2^a\ge8$ stated in print, and record it only because the argument is the level-four one verbatim.

\begin{proposition}[the level invariant; Murty and Thangadurai]\label{prop:trace}
Let $K_{2^a}=\frac2p\sum_{n\in H_{2^a}}\langle n\rangle-|H_{2^a}|$. Then
\[
K_{2^a}=d_0=2^{1-a}\Tr_{\Q(\zeta_{2^a})/\Q}B_{1,\chi}.
\]
In particular $K_4=k$, and $K_{2^a}$ is an odd integer when $p$ is not a Fermat prime.
\end{proposition}

\begin{proof}
The conjugates of $B_{1,\chi}$ are the $B_{1,\chi^u}$ for odd $u$ modulo $2^a$, so the trace is
$\frac1p\sum_n\langle n\rangle\sum_{u\ \mathrm{odd}}\chi(n)^u$. For $z=\chi(n)$ a $2^a$-th root of unity,
$\sum_{w<2^{a-1}}z^{2w+1}=z\sum_w(z^2)^w$ is $2^{a-1}$ if $z=1$, $-2^{a-1}$ if $z=-1$, and $0$ otherwise. Hence the trace is
$2^{a-1}(S_0-S_{2^{a-1}})=2^{a-1}d_0$, and $d_0=2S_0-|H_{2^a}|=K_{2^a}$ by the pairing.
\end{proof}
\leanline{\lean{level_trace}, in the form $\sum_{w<2^{a-1}}\sum_n\chi^{2w+1}(n)\langle n\rangle=2^{a-1}(pS_0-pS_{2^{a-1}})$; the
identification of the sum over odd exponents with the field trace is the classical description of the Galois
conjugates of $B_{1,\chi}$.}

This is Theorem~1 of Murty and Thangadurai \cite{MT}, there read as a statement about the digits of $1/p$, and restated at
the $2$-power levels by Girstmair \cite[Thm.~1]{GirMean}. At $a=2$ it is the trace in \eqref{eq:trnorm}; at level two, for
$p\equiv3\pmod4$ and $p>3$, the same coset sum is $K_2=B_{1,\leg\cdot p}=-h(-p)$.

\begin{example}[levels eight]\label{ex:eight}
At $p=41$ ($a=3$), $H_8=\{1,10,16,18,37\}$ has sum $82=2\cdot41$, so $K_8=2\cdot2-5=-1$. At $p=73$ ($a=3$),
$H_8=\{1,2,4,8,16,32,37,55,64\}$ has sum $219=3\cdot73$, so $K_8=6-9=-3$; with the primitive root $g=5$ and
$\chi(5)=\zeta_8$, the coordinates are $(d_0,d_1,d_2,d_3)=(-3,-3,3,1)$, all odd.
\end{example}

Beyond these facts the higher levels enter the eighth- and sixteenth-range sums through \eqref{eq:suz} at $r=8$ and
$r=16$, together with Bernoulli numbers of twists of conductor $8p$ and $16p$ \cite[\S10]{Berndt}, \cite{GirMean}; as at
level four, these linear formulas do not order the conjugates.

\clearpage
\section{Attribution register}\label{sec:register}

The register records the classical inputs used above.

\begin{center}\small
\begin{tabular}{p{0.5\textwidth}p{0.42\textwidth}}
\toprule
statement & source\\
\midrule
quadratic excess on $(0,p/2)$ and $(0,p/4)$ & Dirichlet 1839; see \cite[(1.1), (1.2)]{Berndt}\\
interval sums as linear combinations of $L(1,\chi\psi)$, every modulus and moment & \cite[(6)]{SUZ}, after Lerch \cite{Lerch},
Berndt \cite{Berndt}, Urbanowicz \cite{Urb}; conjugated form as in \cite[(27)]{GirMean}\\
eighth- and sixteenth-range character sums & \cite[\S7, \S10]{Berndt}\\
moment-weighted interval sums & \cite[\S13]{Berndt}; equation \eqref{eq:suz} at $m\ge2$\\
relative class number as a group determinant in coset sums, through degree ten & \cite{GirMC}, restated in \cite{Kucera}\\
relative class number of the prime-conductor quartic field in coset sums & Hudson--Williams \cite{HW81}; proof in Hudson
\cite[\S2]{Hudson}\\
quarter-count / coset-sum conversion in the quartic case & Hudson \cite[Thm.~3.1]{Hudson}\\
quartic Jacobi-sum normalization $a^2+b^2=p$, $a\equiv-1\pmod4$ & \cite[Cor.~3.10]{BEW}\\
sign of $b$ selected by the pin: $b\equiv-2\pmod8$, i.e.\ $\chi(2)=i^{ab/2}$ & supplementary law for $2$; \cite[Ch.~9]{IR}\\
coset-sum family and deviation vector; norm-type $L$-value product & \cite{GirAbh}, surveyed in \cite{GirSurv}\\
digits of $1/p$ and $h(-p)$ & \cite{GirDig}; \cite[Satz 11]{GirPer}\\
trace of first generalized Bernoulli numbers in the digit formula & \cite[Thm.~1]{MT}; \cite[Thm.~1]{GirMean}\\
$B_{1,\psi}=0$ for even $\psi\ne1$ & \cite[\S4.1]{Washington}\\
integrality of coset sums of a nontrivial subgroup & \cite{GirAbh}, \cite{GirSurv}\\
parity of the coordinates at levels $2^a\ge8$ & argument of Proposition~\ref{prop:odd}; not found in print\\
fixed-base twisted digit identity \eqref{eq:fixed} & \cite[Thm.~1]{GirSurv}\\
digit frequencies as class numbers and classification of when this is possible & \cite{GirDecOct}; \cite{GirFreq}\\
digit variance and Dedekind sums & \cite{GirVar97}; \cite{GirVar26}; \cite{GirMean}\\
$p\equiv1\pmod4$ digit symmetry via composite discriminant & \cite{CK}\\
quadratic residue counts in arithmetic progressions via the class-number formula & \cite{CRST}\\
absence of bias in the function-field analogue & \cite{Rudnick}\\
\bottomrule
\end{tabular}
\end{center}

\section{Formalization}\label{sec:lean}

The accompanying Lean~4 development \lean{QuarticSurvey} (Lean \texttt{v4.35.0}, Mathlib \texttt{v4.35.0};
\url{https://github.com/jkfarland-bot/KSurvey}, release v1.1, archived on Zenodo, all versions
\href{https://doi.org/10.5281/zenodo.23624362}{doi:10.5281/zenodo.23624362}) works with characters
$\chi:\Z/p\Z\to\Z[i]$ of order four, written \lean{MulChar (ZMod p) GaussianInt}, and, for Section~\ref{sec:levels},
with complex-valued characters \texttt{MulChar (ZMod p)}~$\mathbb C$, and with sums over least nonnegative residues. It contains no
\texttt{sorry} and no \texttt{native\_decide}, and every result below depends only on the axioms \texttt{propext},
\texttt{Classical.choice} and \texttt{Quot.sound}. The hypotheses are those of the text: $p$ prime, $p\equiv5\pmod8$, and
$p>5$ where the coset sums enter. In Section~\ref{sec:levels} the hypotheses are stated without $a$: $\chi(-1)=-1$ for
oddness, and $\chi^m=1$ for some $0<m<p-1$ in place of ``$p$ is not a Fermat prime'' (this gives a nontrivial kernel).

\begin{center}\small
\begin{tabular}{lll}
\toprule
statement & Lean & file\\
\midrule
Prop.~\ref{prop:boundary}, odd orders are even & \lean{apply_neg_one_of_odd_orderOf} & Basic\\
Prop.~\ref{prop:boundary}, order two is real & \lean{apply_of_orderOf_two} & Basic\\
Prop.~\ref{prop:boundary}, order four: odd iff $p\equiv5\pmod8$ & \lean{quartic_apply_neg_one}, & Basic\\
 & \lean{quartic_apply_neg_one_of_one} & \\
$\chi_4(2)=\pm i$ & \lean{quartic_apply_two} & Basic\\
Def.~\ref{def:pin}, the pin exists and is unique & \lean{existsUnique_pinned} & Pin\\
the two quartic characters are $\chi_4$ and $\overline{\chi_4}$ & \lean{quartic_eq_or_eq_pow_three} & Pin\\
Prop.~\ref{prop:coset}, integrality & \lean{p_dvd_cosetSum} & CosetSums\\
Prop.~\ref{prop:coset}, class sizes $(p-1)/4$ & \lean{card_class} & CosetSums\\
Prop.~\ref{prop:coset}, pairing $S_j+S_{j+2}=(p-1)/4$ & \lean{S_add_S_neg} & CosetSums\\
Prop.~\ref{prop:coset}, $\beta=k+it$ & \lean{T_eq} & CosetSums\\
Prop.~\ref{prop:odd}, oddness & \lean{odd_k}, \lean{odd_t} & CosetSums\\
conjugation fixes $k$, negates $t$ & \lean{k_pow_three}, \lean{t_pow_three} & CosetSums\\
Prop.~\ref{prop:data}, trace & \lean{trace_eq} & Reconstruction\\
Prop.~\ref{prop:jacobi}, $a^2+b^2=p$ & \lean{jacobiSum_norm} & Jacobi\\
Prop.~\ref{prop:jacobi}, $a\equiv-1$, $b\equiv2\pmod4$ & \lean{jacobiSum_mod_four} & Jacobi\\
Prop.~\ref{prop:jacobi}, $J(\overline\chi,\overline\chi)=\overline{J(\chi,\chi)}$ & \lean{jacobiSum_pow_three} & Jacobi\\
Prop.~\ref{prop:quarters}, \eqref{eq:ABkt}, \eqref{eq:ktAB}, \eqref{eq:hudson} & \lean{two_A}, \lean{two_B}, \lean{five_k}, \lean{five_t}, & Quarters\\
 & \lean{hudson_S} & \\
Thm.~\ref{thm:recon}, formal part & \lean{reconstruction} & Reconstruction\\
Prop.~\ref{prop:fixed} & \lean{fixed_base_collapse} & FixedBase\\
\eqref{eq:suz} at $m=1$, $N=p$, $r=2$ & \lean{half_sum} & FixedBase\\
Prop.~\ref{prop:even}, character and coset forms & \lean{even_T_eq_zero}, \lean{levelSum_even} & Levels\\
Prop.~\ref{prop:levels}, integrality and pairing & \lean{p_dvd_levelSum}, \lean{levelSum_add_neg} & Levels\\
Prop.~\ref{prop:levels}, parity & \lean{level_parity}, \lean{card_kernel_odd} & Levels\\
Prop.~\ref{prop:trace} & \lean{level_trace} & Levels\\
\bottomrule
\end{tabular}
\end{center}

Quoted and not formalized: the identification $\Nm(\beta)=2h^*(p)$ in Proposition~\ref{prop:data}, which needs relative class
numbers of number fields; the $L$-value normalization of Section~\ref{sec:classical}; the general formula \eqref{eq:suz};
the Remark $b\equiv-2\pmod8$, which needs the quartic character of $2$; the identification of the sum over odd exponents
with the field trace in Proposition~\ref{prop:trace}; and the numerical Examples~\ref{ex:13} and~\ref{ex:eight}, which are
checked by direct computation.

\section{Questions}\label{sec:questions}

These record gaps in the literature as found, not claims.

\begin{question}
Berndt's pairing at the sixteenth range \cite[\S10]{Berndt} and the pairing that appears in coordinate congruences at
level eight both arise when characters of order greater than two enter. Is there a uniform statement of this pairing
across the $2$-power levels?
\end{question}

\begin{question}
Does the quartic-order Bernoulli number in \cite[Rem.~5]{GirDecOct}, of conductor $5p$, admit a coordinate description
analogous to the prime-conductor case: a Galois-invariant coordinate determined before pinning and a complementary
coordinate with a natural arithmetic normalization?
\end{question}

\begin{question}
Berndt \cite[\S14]{Berndt} reports, as a written communication from Elliott, that for characters $\chi_p$ of fixed order
in a residue class the values $\arg L(1,\chi_p)$ are everywhere dense modulo $2\pi$. What is the published form of that
result?
\end{question}

\begin{thebibliography}{99}

\bibitem{Berndt} B.~C.~Berndt, \emph{Classical theorems on quadratic residues}, Enseign.\ Math.\ (2) \textbf{22} (1976),
261--304.

\bibitem{CK} K.~Chakraborty and K.~Krishnamoorthy, \emph{On some symmetries of the base $n$ expansion of $1/m$: the class
number connection}, Pacific J.\ Math.\ \textbf{319} (2022), 39--53.

\bibitem{CRST} J.~Chattopadhyay, B.~Roy, S.~Sarkar and R.~Thangadurai, \emph{Distribution of residues modulo $p$ using the
Dirichlet's class number formula}, arXiv:1810.00227 (2019).

\bibitem{GirMC} K.~Girstmair, \emph{The relative class numbers of imaginary cyclic fields of degrees $4$, $6$, $8$ and
$10$}, Math.\ Comp.\ \textbf{61} (1993), 881--887.

\bibitem{GirDig} K.~Girstmair, \emph{The digits of $1/p$ in connection with class number factors}, Acta Arith.\ \textbf{67}
(1994), 381--386.

\bibitem{GirPer} K.~Girstmair, \emph{Periodische Dezimalbr\"uche -- was nicht jeder dar\"uber wei\ss}, Jahrbuch \"Uberblicke
Mathematik 1995, Vieweg, Braunschweig, 1995, 163--179.

\bibitem{GirAbh} K.~Girstmair, \emph{Class number factors and distribution of residues}, Abh.\ Math.\ Sem.\ Univ.\ Hamburg
\textbf{67} (1997), 65--104.

\bibitem{GirVar97} K.~Girstmair, \emph{Digit variance and Dedekind sums}, J.\ Number Theory \textbf{65} (1997), 197--205.

\bibitem{GirDecOct} K.~Girstmair, \emph{On the decimal and octal digits of $1/p$}, Arch.\ Math.\ (Basel), to appear;
arXiv:2510.07873.

\bibitem{GirMean} K.~Girstmair, \emph{Mean values and variances of the digits of $1/p$}, arXiv:2606.29930 (2026).

\bibitem{GirSurv} K.~Girstmair, \emph{Power residues, digit expansions, and relative class numbers}, Surveys Math.\ Appl.\
\textbf{21} (2026), 157--176.

\bibitem{GirFreq} K.~Girstmair, \emph{Digit frequencies and class numbers}, arXiv:2606.04512 (2026).

\bibitem{GirVar26} K.~Girstmair, \emph{On the variance of the digits of $1/p$}, Ramanujan J.\ \textbf{70} (2026), Art.~27.

\bibitem{HW81} R.~H.~Hudson and K.~S.~Williams, \emph{A class number formula for certain quartic fields}, Carleton
Mathematical Series No.~174, Department of Mathematics and Statistics, Carleton University, Ottawa, 1981.

\bibitem{Hudson} R.~H.~Hudson, \emph{Class numbers of imaginary cyclic quartic fields and related quaternary systems},
Pacific J.\ Math.\ \textbf{115} (1984), 129--142.

\bibitem{Kucera} R.~Ku\v{c}era, \emph{Formulae for the relative class number of an imaginary abelian field in the form of a
determinant}, Nagoya Math.\ J.\ \textbf{163} (2001), 167--191.

\bibitem{Lerch} M.~Lerch, \emph{Essai sur le calcul du nombre de classes de formes quadratiques binaires aux coefficients
entiers}, Acta Math.\ \textbf{29} (1905), 333--424.

\bibitem{MT} M.~R.~Murty and R.~Thangadurai, \emph{The class number of $\Q(\sqrt{-p})$ and digits of $1/p$}, Proc.\ Amer.\
Math.\ Soc.\ \textbf{139} (2011), 1277--1289.

\bibitem{BEW} B.~C.~Berndt, R.~J.~Evans and K.~S.~Williams, \emph{Sums of Gauss, Jacobi, and Jacobsthal}, J.\ Number
Theory \textbf{11} (1979), 349--398.

\bibitem{Rudnick} Z.~Rudnick, \emph{On the sum of digits of $1/M$ in $\F_q[x]$}, arXiv:2108.10142 (2021).

\bibitem{SUZ} J.~Szmidt, J.~Urbanowicz and D.~Zagier, \emph{Congruences among generalized Bernoulli numbers}, Acta Arith.\
\textbf{71} (1995), 273--278.

\bibitem{Urb} J.~Urbanowicz, \emph{Connections between $B_{2,\chi}$ for even quadratic Dirichlet characters $\chi$ and class
numbers of appropriate imaginary quadratic fields I}, Compositio Math.\ \textbf{75} (1990), 247--270.

\bibitem{IR} K.~Ireland and M.~Rosen, \emph{A Classical Introduction to Modern Number Theory}, 2nd ed., Graduate Texts in
Mathematics 84, Springer, New York, 1990.

\bibitem{Washington} L.~C.~Washington, \emph{Introduction to Cyclotomic Fields}, 2nd ed., Graduate Texts in Mathematics 83,
Springer, New York, 1997.

\end{thebibliography}

\end{document}
